\documentclass[11pt]{article}

\usepackage[T1]{fontenc}
\usepackage[utf8]{inputenc}
\usepackage[margin=1in]{geometry}
\usepackage{amsmath,amssymb,amsthm,mathtools}
\usepackage{booktabs}
\usepackage{enumitem}
\usepackage{xcolor}
\usepackage{microtype}
\usepackage{hyperref}

\hypersetup{
  colorlinks=true,
  linkcolor=blue!45!black,
  citecolor=blue!45!black,
  urlcolor=blue!55!black
}

\newtheorem{theorem}{Theorem}[section]
\newtheorem{proposition}[theorem]{Proposition}
\newtheorem{lemma}[theorem]{Lemma}
\newtheorem{corollary}[theorem]{Corollary}
\theoremstyle{definition}
\newtheorem{definition}[theorem]{Definition}
\theoremstyle{remark}
\newtheorem{remark}[theorem]{Remark}

\newcommand{\Z}{\mathbb Z}
\newcommand{\Q}{\mathbb Q}
\newcommand{\F}{\mathbb F}
\newcommand{\Gal}{\operatorname{Gal}}
\newcommand{\Frob}{\operatorname{Frob}}
\newcommand{\ord}{\operatorname{ord}}
\newcommand{\vtwo}{v_2}
\newcommand{\cyc}{\mathrm{cyc}}
\newcommand{\itlaw}{\mathrm{it}}

\title{Realization of Gauge-Reduced 2-Power Residue Matrices\\
\large Exact Chebotarev Laws and an Application to Barker Sequences}
\author{Malek Alhazmi\\CytherAi\\\href{https://cytherai.com}{cytherai.com}}
\date{Preprint, July 2026}

\begin{document}
\maketitle

\begin{abstract}
For a prime $p\equiv1\pmod4$ such that $\ord_p(2)$ is odd, the
$2$-primary quotient of $(\Z/p^2\Z)^\times$ is cyclic of order
$2^{\vtwo(p-1)}$.  Its coordinate is defined only up to multiplication by
an odd unit, giving a natural row gauge on matrices of mutual power-residue
data.  We prove that, after this gauge reduction, quadratic-reciprocity
parity and the conditions involving $2$ are the only realizability
constraints.  More precisely, for every fixed realized prefix and every
prescribed depth $t\ge3$, the next-prime transition is an exact
prefix-independent Chebotarev law: one fair parity bit per new edge,
independent uniform coordinates in the corresponding parity classes, and a
common gauge on the new row.  The exact-depth slice has density
$2/4^t$, and every admissible matrix is realized by infinitely many ordered
tuples of distinct primes.  The proof computes a Kummer radical with its
single exceptional $2$-line, determines the maximal abelian subfield by an
affine commutator calculation, and identifies the resulting fiber product
with quadratic reciprocity.  As an application, the reciprocity skeleton
used in a companion Barker-sequence census is the exact iterated arithmetic
law, not merely a null model.  This settles support and iterated-frequency
questions while leaving the simultaneous-height regime, and all
Barker-specific obstructions, open.
\end{abstract}

\begin{center}
\small
\textbf{MSC 2020:} 11R44, 11A15; secondary 11B83, 94A55.\\
\textbf{Keywords:} power residues, Kummer extensions, Chebotarev density,
quadratic reciprocity, residue matrices, Barker sequences.
\end{center}

\section{Introduction}

Let $p$ be an odd prime and put $t_p=\vtwo(p-1)$.  The reduction map
$(\Z/p^2\Z)^\times\to\F_p^\times$ has kernel of odd order, so it induces an
isomorphism on maximal $2$-power quotients.  We write
\[
  \chi_p:(\Z/p^2\Z)^\times\longrightarrow C_{2^{t_p}}
\]
for the resulting quotient character, in additive notation.  A choice of
generator identifies the target with $\Z/2^{t_p}\Z$; changing the generator
multiplies the entire row by one odd unit.  Thus a matrix
$(\chi_{p_i}(p_j))_{i\ne j}$ has an intrinsic row-gauge class even though no
individual numerical coordinate is canonical.

We consider the family
\[
  \mathcal H=\{p:p\equiv1\pmod4,\ \ord_p(2)\text{ is odd}\}.
\]
The terminology ``hard prime'' is used for this family in the accompanying
Barker-sequence code, but is not standard and will not be used in theorem
statements.  If $p\in\mathcal H$, then $\chi_p(2)=0$ and necessarily
$p\equiv1\pmod8$, so $t_p\ge3$.

Quadratic reciprocity gives the visible compatibility relation
\begin{equation}\label{eq:parity-reciprocity}
  \chi_p(q)\equiv\chi_q(p)\pmod2
  \qquad(p,q\in\mathcal H,\ p\ne q).
\end{equation}
Our main result says that there are no others after row gauge.

\begin{theorem}[Exact extension law]\label{thm:extension}
Fix distinct primes $P=(p_1,\ldots,p_r)$ in $\mathcal H$, with
$t_i=\vtwo(p_i-1)$.  Fix $t\ge3$ and put $N=2^t$.  As $q$ varies over
primes in $\mathcal H$ of exact depth $\vtwo(q-1)=t$, the following
conditional distribution exists and is independent of the numerical
prefix $P$:
\begin{enumerate}[label=(\roman*)]
  \item the parities on the $r$ new edges are independent and fair;
  \item conditional on those parities, the column coordinates
  $\chi_{p_i}(q)\in\Z/2^{t_i}\Z$ are independent and uniform in their
  parity classes;
  \item the row coordinates $\chi_q(p_i)\in\Z/N\Z$ are independently
  uniform in the matching parity classes;
  \item the new row is observed modulo simultaneous multiplication by
  $(\Z/N\Z)^\times$.
\end{enumerate}
The rational-prime law on gauge orbits is therefore the pushforward of the
uniform raw-coordinate law.  The unconditional density of the slice is
\[
  \delta_t=\frac{2}{4^t},
\]
also independent of $P$.
\end{theorem}

\begin{corollary}[Realization]\label{cor:realization}
For every depth vector $(t_1,\ldots,t_k)$ with $t_i\ge3$, every
row-gauge class satisfying \eqref{eq:parity-reciprocity} is realized by
infinitely many ordered tuples of distinct primes in $\mathcal H$.
\end{corollary}

The theorem is stronger than a support result: it supplies an exact
per-prefix transition law.  It is also deliberately narrower than a
simultaneous-height theorem.  All prime bounds may be taken successively,
because the conditional law is prefix-independent; nothing here gives a
uniform Chebotarev error term as the prefix field grows.

Corollaries~\ref{cor:product} and~\ref{cor:rado} repackage the theorem as
a single object: an exact arithmetic random structure on $\mathcal H$.
The quadratic-residue parity alone already realizes the countable Rado
graph on this thin family; the full structure equips that skeleton with
independent geometric depth marks, uniform higher $2$-power coordinates,
and a row gauge --- a gauge-valued $2$-adic lift of the prime
quadratic-residue Rado graph, together with its exact extension law.
The Barker application below is one motif family inside this object.

The classical single-row precedent is Schinzel's extension of Kummer's
theorem on prescribed power residues~\cite{Schinzel75,Schinzel77,Schinzel80}.
His Theorem~4 is stated for general exponent $n$, not only odd-prime
exponent; over $\Q$ its unrestricted clause includes $n=2^t$, since the
number $w$ of roots of unity in $\Q$ is $2\not\equiv0\pmod4$.  Applied to
multiplicatively independent rational primes, it realizes arbitrary
single-row power-residue prescriptions.  What it does not supply is the
coupled problem here, where the prescribed radicands are themselves the
prime moduli, the reciprocal row and column must be realized together,
exact depth and the condition on $2$ select one Frobenius fiber, and the
observable row is defined only up to gauge.

Quadratic residue matrices were characterized by Dummit, Dummit, and
Kisilevsky, who also treated cubic and quartic analogues over cyclotomic
ground fields~\cite{DDK}.  For larger exponents they identify two
complications, not one: higher-power reciprocity itself becomes more
complicated, and prime ideals of the cyclotomic base need not be principal.
Their closing remark also considers dropping the primary normalization and
observes that the resulting larger matrix classes appear less tractable to
characterization.  The present paper answers that latter difficulty for one
specific family: rational-prime radicands at arbitrary $2$-power depth.  The
row-gauge quotient records the discarded generator normalization and restores
an exact characterization and transition law.  Like their cubic and quartic
realization proofs, ours adjoins one prime at a time by Chebotarev; the new
content is the quantitative law for that coupled step.  We make no claim that
this family is equivalent to their general higher-power formulation or that
the general $m>4$ problem is settled.  Dummit later developed the
$d$th-power residue-matrix program in the finite-field direction~\cite{DummitFF};
the present extension instead stays over $\Q$ and moves up the $2$-power tower.

DDK also record that direct finite-height counts of their quadratic
configuration types can remain far from the exact Chebotarev frequencies,
and cite the small-prime bias theorem of Dummit--Granville--Kisilevsky as the
explanation~\cite{DDK,DGK}.  This is the direct precedent for the companion
census's opposite-signed finite-height residuals.  It does not resolve them:
the DGK theorem treats fixed Dirichlet characters of fixed conductor, whereas
the census conditions a nonlinear statistic of mutually varying prime moduli.

The multiquadratic field $H$ also occurs naturally in the theory of
minimally tamely ramified extensions.  Dummit--Kisilevsky show that the
decomposition data of such multiquadratic fields are governed by quadratic
residue matrices~\cite{DK}, in the broader minimal-ramification setting of
Kisilevsky--Sonn~\cite{KS}.  These results explain the role of $H$ but do
not give the higher $2$-power intersection $F\cap M=H$ proved below.

The Kummer intersection used below is a specialization of the explicit
theory of Perucca--Sgobba--Tronto~\cite{PST}.  We retain a self-contained
commutator calculation because it exposes exactly where the exceptional
$2$-line enters.  Their degree uniformity is in the exponents for each fixed
finitely generated group $G$; the constants and finite case distinction may
depend on $G$.  The simultaneous-height problem here requires uniformity on
the different, unclaimed axis where
$G=\langle2,p_1,\ldots,p_r\rangle$ gains generators and the radicands grow.

\section{Gauge-reduced character matrices}

\begin{definition}[Raw and gauge-reduced matrices]
Fix depths $\mathbf t=(t_1,\ldots,t_k)$ and put $N_i=2^{t_i}$.  A raw
matrix is an off-diagonal array
\[
  A=(a_{ij}),\qquad a_{ij}\in\Z/N_i\Z.
\]
It is \emph{admissible} if $a_{ij}\equiv a_{ji}\pmod2$ for every
$i\ne j$.  The gauge group
\[
  \mathcal U_{\mathbf t}=\prod_{i=1}^k(\Z/N_i\Z)^\times
\]
acts rowwise by $(u_i)\cdot a_{ij}=u_i a_{ij}$.  An admissible
gauge-reduced matrix is an orbit under this action.
\end{definition}

The condition is well-defined on orbits because every gauge multiplier is
odd.  It is necessary for prime matrices by quadratic reciprocity: the
parity of $\chi_p(q)$ records the Legendre symbol $(q/p)$, and primes in
$\mathcal H$ are $1\pmod4$.  We use the standard power-residue-symbol
conventions of Lemmermeyer~\cite[Chapter~4]{Lemmermeyer}; the additive
coordinate above records the same cyclic quotient while keeping its
generator ambiguity explicit.

\begin{lemma}[The quotient factors through $p$]\label{lem:factor}
The reduction map $(\Z/p^2\Z)^\times\to\F_p^\times$ induces an isomorphism
on maximal $2$-power quotients.  Hence $\chi_p(q)$ may be read from the
Artin symbol of $q$ in the unique degree-$2^{t_p}$ subfield of
$\Q(\zeta_p)$.
\end{lemma}

\begin{proof}
The kernel has order $p$, which is odd.  It therefore lies in the unique
odd-order subgroup and disappears in the $2$-primary quotient.
\end{proof}

\section{The governing compositum}

Fix a realized prefix $P=(p_1,\ldots,p_r)$ and a new depth $t\ge3$.  Put
$N=2^t$ and
\[
  B=\Q(\zeta_{2N}),\qquad B_0=\Q(\zeta_N),\qquad
  \omega=\zeta_N.
\]
For each $i$, let $F_i\subset\Q(\zeta_{p_i})$ be the unique cyclic
subfield of degree $2^{t_i}$, and set
\[
  F=F_1\cdots F_r,\qquad
  H=\Q(\sqrt{p_1},\ldots,\sqrt{p_r}).
\]
Finally choose $\theta_i=p_i^{1/N}$ for $i\ge1$ and the positive
$\theta_0=2^{1/N}$, so that $\theta_0^{N/2}=\sqrt2$, and set
\[
  M=B(\theta_0,\theta_1,\ldots,\theta_r).
\]
The field $F$ carries the old-row/new-column coordinates; $M$ carries the
new row and the condition involving $2$.

\begin{proposition}[The radical and its $2$-line]\label{prop:radical}
In $B^\times/B^{\times N}$ one has
\[
 \langle2,p_1,\ldots,p_r\rangle B^{\times N}/B^{\times N}
 \cong \Z/(N/2)\Z\oplus(\Z/N\Z)^r.
\]
Consequently $[M:B]=N^r(N/2)$.
\end{proposition}

\begin{proof}
Suppose $2^{e_0}\prod_i p_i^{e_i}=\beta^N$ in $B$.  The field $B$ is
ramified only at $2$.  At a prime of $B$ over $p_j$, the normalized
valuation of $p_j$ is $1$, hence $N\mid e_j$ for every $j$.  Thus no mixed
relation survives and only the class of $2$ remains.

That class has order at most $N/2$, since
\[
  2^{N/2}=(\sqrt2)^N,
  \qquad \sqrt2=\zeta_8+\zeta_8^{-1}\in B.
\]
Its order is not smaller.  Otherwise $2^{N/4}=\beta^N$ for some $\beta$,
so $2^{1/4}=\beta\zeta$ for an $N$th root of unity $\zeta\in B$.  This
would embed the non-Galois quartic field $\Q(2^{1/4})$ in the abelian
extension $B/\Q$, impossible because every subfield of an abelian Galois
extension is Galois.  Thus the order is exactly $N/2$.
\end{proof}

The exponent is divisible by $8$, so the exceptional $2$-line behind the
Grunwald--Wang special case is a mandatory audit rather than a remote
analogy~\cite{Wang}.  The Kummer calculation here is over
$B\supset\mu_N$, so no local--global existence theorem is being invoked.
Instead the possible collapse is computed directly and exhausted: the
displayed relation gives the single order drop on the $2$-line, while the
valuation argument and the lower bound from
$\Q(2^{1/4})\not\subset B$ exclude any second radical relation.  This drop
will appear below as an even translation coordinate.

\section{The maximal abelian subfield}

Write $T=\Gal(M/B)$.  By Proposition~\ref{prop:radical}, an element of $T$
has coordinates $c=(c_0,c_1,\ldots,c_r)$ with
\[
  \tau_c(\theta_i)=\omega^{c_i}\theta_i,
  \qquad c_0\equiv0\pmod2.
\]
Let $G=\Gal(M/\Q)$.  For $a\in(\Z/2N\Z)^\times$, choose a lift
$\widetilde\sigma_a$ of $\sigma_a\in\Gal(B/\Q)$ that fixes the
$\theta_i$ for $i\ge1$ and satisfies
\[
  \widetilde\sigma_a(\theta_0)=\omega^{u(a)}\theta_0,
  \qquad (-1)^{u(a)}=\epsilon(a):=\frac{\sigma_a(\sqrt2)}{\sqrt2}.
\]
Such lifts exist by Proposition~\ref{prop:radical}.  Moreover
\begin{equation}\label{eq:epsilon}
  \epsilon(a)=1\quad\Longleftrightarrow\quad a\equiv\pm1\pmod8.
\end{equation}

\begin{proposition}[Maximal abelian subfield]\label{prop:ab}
One has
\[
  [G,G]=\Gal(M/BH),\qquad M^{\mathrm{ab}}=BH.
\]
\end{proposition}

\begin{proof}
Conjugation gives
$\widetilde\sigma_a\tau_c\widetilde\sigma_a^{-1}=\tau_{ac}$, so
commutators with $T$ generate
\[
  T^2=\{c_0\equiv0\pmod4,\ c_i\equiv0\pmod2\ (i\ge1)\}.
\]
The only missing translation is the second bit on the $2$-line.  Use
left-translation coordinates, with multiplication
\[
  (a,u)(b,v)=(ab,u+av).
\]
\begin{equation}\label{eq:commutator}
 [(a,u),(b,v)]=\bigl(1,(a-1)v-(b-1)u\bigr).
\end{equation}
For $a=3$ and $b=5$, both $u(3)$ and $u(5)$ are odd by
\eqref{eq:epsilon}; hence the translation in \eqref{eq:commutator} is
$2\pmod4$.  Together with $T^2$ this generates
\[
  \{c_0\equiv0\pmod2,\ c_i\equiv0\pmod2\ (i\ge1)\}
  =\Gal(M/BH).
\]
The reverse containment follows because $BH/\Q$ is abelian.
\end{proof}

There is an independent external check.  Apply
Perucca--Sgobba--Tronto~\cite[Theorem~5.1]{PST} to the positive,
torsion-free group $\langle2,p_1,\ldots,p_r\rangle$, with their parameters
$m=t+1$ and $n=t$.  Their formula gives
\[
  M\cap\Q^{\cyc}=B(\sqrt2,\sqrt{p_1},\ldots,\sqrt{p_r})=BH.
\]
Kronecker--Weber then gives $M^{\mathrm{ab}}=BH$, agreeing with
Proposition~\ref{prop:ab}.

\begin{corollary}[The only intersection is reciprocity]\label{cor:intersection}
\[
  F\cap M=H.
\]
\end{corollary}

\begin{proof}
Since $F/\Q$ is abelian, $F\cap M\subseteq M^{\mathrm{ab}}=BH$.  The field
$B$ is totally ramified at $2$, whereas $F$ and $H$ are unramified at
$2$; consequently the inertia-fixed field of $BH$ at $2$ is $H$ and
$F\cap BH\subseteq H$.  Conversely, $\sqrt{p_i}$ lies in both $F_i$ and
$M$, so $H\subseteq F\cap M$.
\end{proof}

The $F_i$ are linearly disjoint over $\Q$: each is ramified only at its
own prime $p_i$, so an intersection with the compositum of the others would
be unramified at every finite prime and is therefore $\Q$.  Hence
\begin{equation}\label{eq:Fdegree}
  \Gal(F/\Q)\cong\prod_{i=1}^r\Z/2^{t_i}\Z.
\end{equation}

\section{Frobenius coordinates and the row gauge}

By Corollary~\ref{cor:intersection}, the usual fiber-product description is
\begin{equation}\label{eq:fiberproduct}
 \Gal(FM/\Q)\cong
 \{(\phi,\psi)\in\Gal(F/\Q)\times\Gal(M/\Q):
   \phi|_H=\psi|_H\}.
\end{equation}
Use Lemma~\ref{lem:factor} to choose the $i$th coordinate in
\eqref{eq:Fdegree} so that the Artin symbol of $q$ is
$c_i=\chi_{p_i}(q)$.  The restriction of this coordinate to
$\Q(\sqrt{p_i})$ is $(-1)^{c_i}$.

The reverse coordinates must be read over $B_0$, not over $B$.  If
$\vtwo(q-1)=t$, then
\[
 q\equiv1+N\pmod{2N}.
\]
Thus $q$ splits completely in $B_0$ but has residue degree $2$ in $B$.
The extension $M/B_0$ is Galois: $B/B_0$ is Galois and $B_0$ contains
$\mu_N$, so every Kummer line used to define $M$ is stable.  Choose
$\mathfrak q_0\mid q$ in $B_0$ and a prime above it in $M$.  Since
the relative Frobenius fixes $\omega$, write
\[
  \Frob_{\mathfrak q_0}(\theta_i)=\omega^{d_i}\theta_i.
\]
Reduction modulo the chosen prime gives
\[
  \omega^{d_i}\equiv p_i^{(q-1)/N}\pmod q.
\]
If $g$ is a primitive root modulo $q$, then
$p_i^{(q-1)/N}=(g^{(q-1)/N})^{\chi_q(p_i)}$.  The reduction of $\omega$
is another generator of the same order-$N$ group.  Therefore
\begin{equation}\label{eq:rowgauge}
  (d_1,\ldots,d_r)=u\,(\chi_q(p_1),\ldots,\chi_q(p_r))
\end{equation}
for one odd unit $u$ common to the row.

Changing the prime of $B_0$ above the rational prime $q$ changes this unit.
Changing only the prime of $M$ above a fixed $\mathfrak q_0$ does not:
an element above $\sigma_{1+N}$ centralizes $T$ because $1+N\equiv1\pmod
N$.  Thus the ambiguity in \eqref{eq:rowgauge} is exactly the row gauge,
not an additional quotient.

Finally, the condition in \eqref{eq:fiberproduct} is precisely
\begin{equation}\label{eq:fiberparity}
  c_i\equiv d_i\pmod2\qquad(1\le i\le r),
\end{equation}
which is quadratic reciprocity and nothing more.

\section{Proof of the exact extension law}

Let $A=\prod_i2^{t_i}=[F:\Q]$.  From the preceding degree calculations,
\begin{equation}\label{eq:grouporder}
 [FM:\Q]=\frac{[F:\Q][M:\Q]}{[H:\Q]}
 =\frac{A\,N^{r+2}}{2^{r+1}}.
\end{equation}
Consider the subset $X_{P,t}\subseteq\Gal(FM/\Q)$ defined by
\[
 \psi|_B=\sigma_{1+N},\qquad \psi(\theta_0)=\theta_0.
\]
This is the exact-depth, odd-$2$-order Frobenius slice.

\begin{lemma}[Every admissible fiber is nonempty]\label{lem:nonempty}
Given arbitrary $c_i\in\Z/2^{t_i}\Z$ and $d_i\in\Z/N\Z$ satisfying
\eqref{eq:fiberparity}, there is a unique element of $X_{P,t}$ having those
coordinates.
\end{lemma}

\begin{proof}
Choose $\phi\in\Gal(F/\Q)$ with coordinates $c_i$.  Define $\psi$ by
\[
 \psi|_B=\sigma_{1+N},\qquad
 \psi(\theta_0)=\theta_0,\qquad
 \psi(\theta_i)=\omega^{d_i}\theta_i.
\]
Proposition~\ref{prop:radical} says that the displayed radical relations
are complete.  The only nontrivial compatibility is
$\theta_0^{N/2}=\sqrt2$, and it holds because $t\ge3$ gives
$1+N\equiv1\pmod8$, so $\sigma_{1+N}(\sqrt2)=\sqrt2$.  Thus $\psi$ is an
automorphism.  On $\sqrt{p_i}$, it acts by $(-1)^{d_i}$; hence
\eqref{eq:fiberparity} is exactly the condition that $(\phi,\psi)$ belongs
to the fiber product \eqref{eq:fiberproduct}.  Uniqueness is immediate
from the generators.
\end{proof}

There are $A$ choices of $c=(c_i)$ and, for each one, $(N/2)^r$ choices of
$d=(d_i)$ with matching parity.  Therefore
\begin{equation}\label{eq:slicesize}
  |X_{P,t}|=A(N/2)^r.
\end{equation}
The slice is stable under conjugacy.  Indeed, the $F$ component and the
restriction to $B$ are abelian, and conjugating an element above
$\sigma_{1+N}$ rescales $d$ by one common unit while keeping the
$\theta_0$ coordinate zero.  In affine notation the possible translation
term is multiplied by $1-(1+N)$, which vanishes modulo $N$.

Chebotarev applied to the union $X_{P,t}$ now gives~\cite{Neukirch}
\[
 \delta_t=\frac{|X_{P,t}|}{[FM:\Q]}
 =\frac{2}{N^2}=\frac{2}{4^t}.
\]
More generally, each raw coordinate vector occurs once in
$X_{P,t}$, while conjugacy identifies precisely the common row-gauge
orbits.  Conditional on the slice, Chebotarev therefore gives the
pushforward of the uniform raw-coordinate distribution.  This proves
Theorem~\ref{thm:extension}.

For $t=2$, the two requirements are incompatible:
$1+N=5\pmod8$ sends $\sqrt2$ to $-\sqrt2$ and cannot fix $\theta_0$.
This proves the depth floor rather than assuming it.  Summing the exact
slices gives
\begin{equation}\label{eq:densitysum}
 \sum_{t\ge3}\frac{2}{4^t}=\frac1{24}.
\end{equation}
The interchange of density and sum is harmless: the upper density of
primes with $\vtwo(p-1)\ge T$ is at most the density
$2^{1-T}$ of $p\equiv1\pmod{2^T}$, which tends to zero.  Equation
\eqref{eq:densitysum} agrees with Hasse's density theorem~\cite{Hasse}:
the odd-order family has density $7/24$, and its $p\equiv7\pmod8$ part has
density $1/4$, leaving $1/24$ in the $p\equiv1\pmod4$ sector.

\begin{proof}[Proof of Corollary~\ref{cor:realization}]
Start with the empty prefix and adjoin the vertices in any chosen order.
For a target gauge class, choose a raw representative.  If the realized
prefix differs by old-row gauges, rescale the corresponding new column
coordinates; the transition law is equivariant because odd multiplication
preserves each parity class.  Lemma~\ref{lem:nonempty} gives a nonempty
conjugacy class for every admissible extension, and Chebotarev supplies
infinitely many new primes in it.  Such primes are unramified in the
compositum and hence distinct from $2,p_1,\ldots,p_r$.  Induction completes
the realization.
\end{proof}

\section{The exact iterated matrix law}

The extension theorem canonically determines a projectively consistent,
exchangeable family of finite probability spaces
(Corollary~\ref{cor:product} below).  For a fixed depth vector the
construction is as follows.  For $i<j$, put
\[
 E_{ij}=\{(a,b)\in\Z/2^{t_i}\Z\times\Z/2^{t_j}\Z:
                    a\equiv b\pmod2\},
\]
and define the raw product measure
\begin{equation}\label{eq:rawmeasure}
 \nu^{\mathrm{raw}}_{\mathbf t}
   =\bigotimes_{i<j}\operatorname{Unif}(E_{ij}).
\end{equation}
Let $\nu_{\mathbf t}$ be its pushforward to row-gauge classes.

\begin{theorem}[Iterated law]\label{thm:iterated}
Fix an ordered depth vector $\mathbf t$.  Draw primes of these depths from
$\mathcal H$ and take their bounds to infinity successively, with the last
adjoined prime innermost.  The resulting distribution of gauge-reduced
character matrices is exactly $\nu_{\mathbf t}$.  It is independent of the
adjoining order.
\end{theorem}

\begin{proof}
For a fixed numerical prefix, Theorem~\ref{thm:extension} supplies the next
factor
\[
  \prod_{i<j}\operatorname{Unif}(E_{ij}),
\]
independently of that prefix.  Induction therefore gives
\eqref{eq:rawmeasure}.  The kernel is equivariant under every old-row gauge,
which rescales the corresponding new column, and under the new-row gauge,
which rescales all new-row entries together.  Hence choosing a raw
representative at an intermediate stage does not affect the final
pushforward.  Permuting the adjoining order merely permutes the factors in
\eqref{eq:rawmeasure}.
\end{proof}

This theorem concerns nested limits.  It does not assert convergence when
all primes are bounded by one common height: the field $FM$ and its
discriminant grow with the prefix, and the ordinary Chebotarev theorem used
above supplies no uniform error term for that regime.

Two corollaries repackage the depth densities and the iterated law as one
unconditional object.

\begin{corollary}[Unconditional marked product law]\label{cor:product}
Within $\mathcal H$ the exact-depth densities of
Theorem~\ref{thm:extension} normalize by \eqref{eq:densitysum} to
\[
  \Pr(T=t)=\frac{2/4^{t}}{1/24}=\frac{3}{4^{t-2}},\qquad t\ge3,
\]
a geometric law summing to $1$.  Draw $k$ primes from $\mathcal H$ by
successive bounds as in Theorem~\ref{thm:iterated}, without prescribing
depths.  Then the depth marks $T_1,\ldots,T_k$ are independent with the
law above; conditional on $\mathbf T=\mathbf t$ the raw matrix is
distributed as $\nu^{\mathrm{raw}}_{\mathbf t}$; explicitly,
\[
 \Pr(\mathbf T=\mathbf t,\ A=a)
 =\prod_{i=1}^{k}\frac{3}{4^{t_i-2}}
  \prod_{1\le i<j\le k}
  \frac{\mathbf 1_{\{a_{ij}\equiv a_{ji}\ (\mathrm{mod}\ 2)\}}}
       {2^{t_i+t_j-1}},
\]
and the gauge-reduced structure follows the row-gauge pushforward.  The
family over $k$ is exchangeable and projectively consistent: permuting
vertices permutes the factors, and deleting a vertex deletes its depth
factor and its incident pair factors.  It therefore extends to an
exchangeable law on countable gauge-reduced structures.  As in
Theorem~\ref{thm:iterated}, every statement is a nested-limit statement;
nothing here concerns simultaneous height.
\end{corollary}

\begin{proof}
Depth independence is the prefix-independence of the slice density in
Theorem~\ref{thm:extension}, normalized by \eqref{eq:densitysum}; the
conditional coordinate law is Theorem~\ref{thm:iterated}, and the display
multiplies the factors $|E_{ij}|^{-1}=2^{-(t_i+t_j-1)}$.  Consistency
holds in raw coordinates by the product form, and passes to gauge classes
because the odd-unit action on a row is coordinatewise, so restriction to
a subset of vertices commutes with the quotient.  The Kolmogorov
extension theorem then applies.
\end{proof}

\begin{corollary}[Rado skeleton]\label{cor:rado}
Join $p,q\in\mathcal H$ when $\chi_p(q)$ is even --- equivalently, by
quadratic reciprocity, when $p$ and $q$ are quadratic residues of one
another.  The resulting countable graph is isomorphic to the Rado graph.
\end{corollary}

\begin{proof}
Let $U,V\subset\mathcal H$ be finite and disjoint.  Over the realized
prefix $U\cup V$, at any depth $t\ge3$, Theorem~\ref{thm:extension} makes
the new parity bits independent and fair, so the pattern ``even toward
every vertex of $U$, odd toward every vertex of $V$'' is a cylinder of
positive density: infinitely many $q\in\mathcal H$ are adjacent to all of
$U$ and to none of $V$.  A countably infinite graph with this extension
property is isomorphic to the Rado graph by
back-and-forth~\cite{Cameron}.
\end{proof}

The analogous construction on all primes $p\equiv1\pmod4$ is a classical
realization of the Rado graph~\cite{Cameron}; the content here is that
the extension property survives restriction to the density-$1/24$
odd-order family, and that this graph is only the parity skeleton of the
full gauge-valued structure.

\begin{remark}[Toward a universal homogeneous gauge structure]
\label{rem:fraisse}
Corollary~\ref{cor:realization} proves the one-point extension property
in the full gauge-valued signature: every admissible extension of every
realized finite substructure is realized by infinitely many primes.
Formalizing the class of finite gauge-reduced matrices with its notion of
embedding, and identifying the arithmetic structure on $\mathcal H$ as
the universal homogeneous object of that class --- isomorphic to almost
every sample of the law of Corollary~\ref{cor:product} --- is left to
subsequent work and is not part of the frozen claim set.
\end{remark}

\section{Application to the Barker reciprocity skeleton}

Turyn's character-sum framework supplies classical necessary conditions
for cyclic difference sets and the Barker problem~\cite{Turyn}.  The
companion census~\cite{Census} studies gauge-invariant predicates of
five-vertex character matrices.  A pair $\{i,j\}$ is witnessed by
$\ell\notin\{i,j\}$ when $a_{\ell i}+a_{\ell j}=0$.  The relevant marked
cell imposes global covering and true minimality, no zero off-diagonal
entries, an all-odd cofactor row at a marked depth-$3$ vertex $x$, and
exactly two pairs witnessed by $x$.  Its success event is
\[
  \sigma_x=\sum_{j\ne x}a_{xj}=0\pmod8.
\]
Both the cell and success are gauge-invariant.

For a full depth vector $\mathbf t$, let $A(M,x)$ denote that eligibility
event and define the marked-cell intensity ratio
\begin{equation}\label{eq:Rit}
 R^{\itlaw}_{\mathbf t}=
 \frac{\sum_{x:t_x=3}\nu_{\mathbf t}(A(M,x)\wedge\sigma_x=0)}
      {\sum_{x:t_x=3}\nu_{\mathbf t}(A(M,x))}.
\end{equation}
The numerator and denominator pool marked cells before division; this is
not the expectation of a per-matrix ratio.

An exhaustive exact contraction of \eqref{eq:rawmeasure} in the companion
repository gives~\cite{Census}
\begin{align}
 R^{\itlaw}_{(3,3,3,3,3)}&=\frac{1373}{5300}=0.2590566\ldots,\\
 R^{\itlaw}_{(3,3,3,3,4)}&=\frac{1123}{4215}=0.2664294\ldots.
\end{align}
The companion paper uses the four-entry cofactor shorthand
$R(3,3,3,3)$ and $R(3,3,3,4)$ for these two quantities.  Theorem
\ref{thm:iterated} upgrades their interpretation: they are exact iterated
arithmetic probabilities, not fitted or Monte Carlo null-model constants.

\begin{corollary}[Scoped Barker consequence]\label{cor:barker}
No skeleton-admissible pattern is globally absent from tuples of primes in
$\mathcal H$.  Therefore a proposed Barker argument cannot discard such a
pattern merely by asserting that its gauge-reduced $2$-power residue matrix
is arithmetically unrealizable.
\end{corollary}

This is a realizability statement only.  A realized prime tuple need not be
the prime support of a Barker candidate, so a Barker-specific necessary
condition may still cut the realized set.  Nor does the theorem address
nonabelian, class-group, or higher-cohomological data not determined by the
matrix.  Finally, it does not explain why a skeleton-admissible state may be
absent below a finite bound, or why a simultaneous-height census rate may
differ from \eqref{eq:Rit}.  Those are frequency questions outside the
support theorem.

\section{Computational guardrail and freeze boundary}

The proof is algebraic.  A deterministic program accompanies it only as a
falsification harness.  It checks the affine commutator through depths
$2$--$6$, the predicted derived subgroup, the exact slice cardinality for
multiple prefix ranks, conjugacy as common row scaling, and $56$
non-vacuous identifications between the repository's $\chi$ coordinates
and the Kummer residue in \eqref{eq:rowgauge}.  A sieve over all
$1{,}270{,}607$ primes below $2\cdot10^7$ gives the exact-depth counts
\[
\begin{array}{c|rrrrr}
t&3&4&5&6&7\\ \hline
\#&39711&9939&2536&619&156
\end{array}
\]
and zero at $t=2$, totaling $53008$ primes in the $1\pmod4$ sector.  These
values reproduce the theorem's density slices numerically; they are not
used as evidence for the proof.

The reproducibility layer is also part of the methodological comparison.
The benchmark Barker computations of Mossinghoff and of
Borwein--Mossinghoff publish search data and candidate lists and describe
their algorithms~\cite{BM1,BM2}.  They do not describe an executable
end-to-end reproduction driver or an independent implementation of the
arithmetic primitives.  The 2014 sequel itself reports that an earlier
implementation of Tarjan's circuit-enumeration algorithm~\cite{Tarjan}
missed cycles, and that earlier software for Turyn's test could not resolve
some cases~\cite{BM2}.  The fail-closed driver, pinned artifacts, regression
suite, and clean-room arithmetic path accompanying the census are designed
around that documented failure mode.  This is a reproducibility advantage,
not a claim of priority over the earlier computations.

The program is
\texttt{barker\_k6\_bundle/research/realization\_checks.py}; its stable output
is \texttt{\_realization\_checks.json}.  The publication freeze excludes
three attractive but logically separate projects:
\begin{enumerate}
  \item a simultaneous-height convergence or bias theorem;
  \item larger Barker censuses or targeted motif enumeration;
  \item combinatorial mechanisms for eligibility-weighting effects in the
  exact skeleton.
\end{enumerate}
Each requires new analytic theory or a new scalable algorithm.  None can
alter Theorems~\ref{thm:extension} and~\ref{thm:iterated}, and none is
needed for Corollary~\ref{cor:barker} at its stated scope.

\section*{Acknowledgment of scope}

The central intersection formula is independently covered by the explicit
Kummer theory of Perucca--Sgobba--Tronto~\cite{PST}.  The contribution here
is the gauge-reduced residue-matrix formulation, the exact Frobenius slice
and transition law, the iterated product law, and their scoped application
to the Barker reciprocity skeleton.  The paper makes no claim to solve the
even Barker conjecture.

\begin{thebibliography}{99}

\bibitem{BM1}
M.~J. Mossinghoff,
\emph{Wieferich pairs and Barker sequences},
Des. Codes Cryptogr. \textbf{53} (2009), no.~3, 149--163,
\href{https://doi.org/10.1007/s10623-009-9301-3}{doi:10.1007/s10623-009-9301-3}.

\bibitem{BM2}
P.~Borwein and M.~J. Mossinghoff,
\emph{Wieferich pairs and Barker sequences, II},
LMS J. Comput. Math. \textbf{17} (2014), no.~1, 24--32,
\href{https://doi.org/10.1112/S1461157013000223}{doi:10.1112/S1461157013000223}.

\bibitem{Cameron}
P.~J. Cameron,
\emph{The random graph},
in: The Mathematics of Paul Erd\H{o}s II
(R.~L. Graham and J.~Ne\v{s}et\v{r}il, eds.),
Algorithms Combin. \textbf{14}, Springer, Berlin, 1997, pp.~333--351.

\bibitem{Census}
M.~Alhazmi,
\emph{Minimal covering configurations of hard primes: a cofactor-cycle
theorem and the arithmetic floor of a 2-primary census}, preprint and
reproducibility archive, 2026.

\bibitem{DDK}
D.~S. Dummit, E.~P. Dummit, and H.~Kisilevsky,
\emph{Characterizations of quadratic, cubic, and quartic residue matrices},
J. Number Theory \textbf{168} (2016), 167--179,
\href{https://doi.org/10.1016/j.jnt.2016.04.014}{doi:10.1016/j.jnt.2016.04.014}.

\bibitem{DK}
D.~S. Dummit and H.~Kisilevsky,
\emph{Decomposition types in minimally tamely ramified extensions of
$\Q$}, Res. Number Theory \textbf{3} (2017), article 24,
\href{https://doi.org/10.1007/s40993-017-0085-7}{doi:10.1007/s40993-017-0085-7}.

\bibitem{DummitFF}
E.~P. Dummit,
\emph{Characterizations of the $d$th-power residue matrices over finite
fields}, arXiv:1809.10225 (2018),
\href{https://arxiv.org/abs/1809.10225}{arXiv:1809.10225}.

\bibitem{DGK}
D.~S. Dummit, A.~Granville, and H.~Kisilevsky,
\emph{Big biases amongst products of two primes},
Mathematika \textbf{62} (2016), no.~2, 502--507,
\href{https://doi.org/10.1112/S0025579315000339}{doi:10.1112/S0025579315000339}.

\bibitem{Hasse}
H.~Hasse,
\emph{\"Uber die Dichte der Primzahlen $p$, f\"ur die eine vorgegebene
ganzrationale Zahl $a\ne0$ von gerader bzw. ungerader Ordnung mod. $p$ ist},
Math. Ann. \textbf{166} (1966), 19--23.

\bibitem{KS}
H.~Kisilevsky and J.~Sonn,
\emph{On the minimal ramification problem for $\ell$-groups},
Compos. Math. \textbf{146} (2010), no.~3, 599--606.

\bibitem{Lemmermeyer}
F.~Lemmermeyer,
\emph{Reciprocity Laws: From Euler to Eisenstein}, Springer Monographs in
Mathematics, Springer, 2000.

\bibitem{Neukirch}
J.~Neukirch,
\emph{Algebraic Number Theory}, Grundlehren der mathematischen
Wissenschaften 322, Springer, 1999.

\bibitem{PST}
A.~Perucca, P.~Sgobba, and S.~Tronto,
\emph{Explicit Kummer theory for the rational numbers},
Int. J. Number Theory \textbf{16} (2020), no.~10, 2213--2231,
\href{https://doi.org/10.1142/S1793042120501146}{doi:10.1142/S1793042120501146}.

\bibitem{Schinzel75}
A.~Schinzel,
\emph{On power residues and exponential congruences},
Acta Arith. \textbf{27} (1975), 397--420.

\bibitem{Schinzel77}
A.~Schinzel,
\emph{Abelian binomials, power residues and exponential congruences},
Acta Arith. \textbf{32} (1977), 245--274,
\href{https://doi.org/10.4064/aa-32-3-245-274}{doi:10.4064/aa-32-3-245-274}.

\bibitem{Schinzel80}
A.~Schinzel,
\emph{Addendum and corrigendum to ``Abelian binomials, power residues and
exponential congruences''}, Acta Arith. \textbf{36} (1980), 101--104,
\href{https://doi.org/10.4064/aa-36-1-101-104}{doi:10.4064/aa-36-1-101-104}.

\bibitem{Tarjan}
R.~Tarjan,
\emph{Enumeration of the elementary circuits of a directed graph},
SIAM J. Comput. \textbf{2} (1973), no.~3, 211--216,
\href{https://doi.org/10.1137/0202017}{doi:10.1137/0202017}.

\bibitem{Turyn}
R.~J. Turyn,
\emph{Character sums and difference sets},
Pacific J. Math. \textbf{15} (1965), 319--346.

\bibitem{Wang}
S.~Wang,
\emph{On Grunwald's theorem},
Ann. of Math. (2) \textbf{51} (1950), 471--484,
\href{https://doi.org/10.2307/1969335}{doi:10.2307/1969335}.

\end{thebibliography}

\end{document}
